% Part: history
% Chapter: biographies
% Section: ernst-zermelo

\documentclass[../../../include/open-logic-section]{subfiles}

\begin{document}

\olfileid{his}{bio}{zer}

\olsection{إرنست زيرميلو}

\olphoto{zermelo-ernst}{Ernst Zermelo}

وُلد إرنست زيرميلو ببرلين في ألمانيا يوم \OLClassicalDigits{27} يوليو \OLClassicalDigits{1871}، وله خمس أخوات.
وكان المرض كثيراً في أسرته، فلم تدرك الرشد من أخواته إلا ثلاث، ومات أبواه
وهو صغير، فصار هو وإخوته يتامى وله سبع عشرة سنة. وكان شديد الميل إلى
الفنون، ولا سيما الشعر، معروفاً بنفاذ ذهنه وسرعة خاطره ونزعته إلى النقد.
ومن أشهر ما خلّفه في الرياضيات إدخاله مسلّمة الاختيار سنة \OLClassicalDigits{1904}، ووضعه سنة
\OLClassicalDigits{1908} صياغة مسلّمية لنظرية المجموعات.

وتشعب علم زيرميلو في الجامعة، فدرس الفيزياء والرياضيات والفلسفة. وأتم في
جامعة برلين سنة \OLClassicalDigits{1894}، تحت إشراف هيرمان شفارتس، أطروحته~\emph{بحوث في
حساب التغيرات}. ثم آثر سنة \OLClassicalDigits{1897} أن يتم درسه في جامعة غوتنغن، فاشتد تأثره
هناك بعمل دافيد هيلبرت في أسس الرياضيات. وتأهل سنة \OLClassicalDigits{1899} لرتبة الأستاذ،
غير أنه لم يُولَّ منصباً إلا بعد إحدى عشرة سنة؛ ولعل ما وُصف به من غرابة
السلوك و«العجلة العصبية» كان سبباً في ذلك.

وآل أمره سنة \OLClassicalDigits{1910} إلى أستاذية ذات راتب في جامعة زيورخ؛ ثم أكرهه السل على
التقاعد سنة \OLClassicalDigits{1916}. فلما برئ منه مُنح سنة \OLClassicalDigits{1921} أستاذية فخرية في جامعة
فرايبورغ، وعكف في تلك المدة على أسس الرياضيات. وقد ساءته آثار تورالف سكولم
وكورت غودل، فجهر في مقالاته بانتقاد طريقتيهما. ثم عُزل سنة \OLClassicalDigits{1935} من منصبه
بفرايبورغ، لقلة حظوته ولمعارضته صعود هتلر إلى الحكم في ألمانيا.

وأمضى زيرميلو أواخر عمره في عزلة؛ فإنه لما عُزل سنة \OLClassicalDigits{1935} ترك الرياضيات
وانتقل إلى الريف، فعاش عيشاً مقتصداً. وتزوج سنة \OLClassicalDigits{1944}، فلما أخذ بصره يضعف
صار عيالاً على زوجه في كل شأن، ثم ذهب بصره كله سنة \OLClassicalDigits{1951}. وتوفي بغونترستال
في ألمانيا يوم \OLClassicalDigits{21} مايو \OLClassicalDigits{1953}.

\begin{reading}
ولسيرة زيرميلو المستوفاة يُنظر~\citet{Ebbinghaus2015}. ومقالتاه المؤسستان
المنشورتان سنتي \OLClassicalDigits{1904} و\OLClassicalDigits{1908} مقروءتان بأصلهما الألماني
\citep{Zermelo1904,Zermelo1908}. ومجموع أعماله، وفيها ما كتبه في
الفيزياء، متاح بترجمة إنجليزية في
\citep{Ebbinghaus2010,Ebbinghaus2013}.
\end{reading}

\end{document}
